\subsection{Coupled-scene prediction protocol}
\label{sec:coupled-protocol}
We separate the material of an object from the interaction applied to it.
A world $w$ specifies fixed object geometry, density, and a hidden material
field $m_w$. An interaction $a$ specifies the initial state, prescribed fixture
positions and release schedule, gravity, and collision geometry. For each
world, a short calibration interaction provides point observations
$\mathcal D_w=\{(t_i,X_j,x_{ij})\}$ up to a fixed cutoff. A method fits its
material representation to this prefix, seals the fitted parameters, and
predicts two different interactions without refitting. The same hidden
material is retained across calibration and both queries.

The protocol contains six scene families: moving and released compliant
clamps; pushing around a curved obstacle; gravity drops into a curved
container, including an obstacle; a translating finite-mass rigid striker;
two deformable objects in collision or stacking; and spatially heterogeneous
material. Query changes include fixture location, displacement amplitude,
impact velocity, drop height and orientation, and contact configuration.
The frozen census is four independently sampled material worlds per
family and two queries per world. Object shape and density are fixed in
this study; transfer is across material and interaction, rather than unseen
object geometry.
The protocol, implementation inventory, model checkpoints, and development
decision were fixed before generating the held worlds. Young's modulus is
sampled log-uniformly from $8$ to $18$\,kPa and Poisson's ratio uniformly
from $0.25$ to $0.40$. Two-body scenes use independent material draws.
Heterogeneous worlds additionally sample a log-contrast corresponding to
factors between two and three and a transition coordinate between $0.35$
and $0.65$. These draws remain fixed across each world's three interactions.

All methods receive the same calibration tracks, rest geometry, density,
and public boundary parameters. Tracks cover a fixed two-face material-point
subset and are sampled every $2$\,ms. This controlled track-input study does
not measure image tracking errors or camera-dependent occlusion. Materials
are unknown, while fixture compliance and collision geometry are known.
Consequently, this experiment evaluates prediction conditional on supplied
boundaries; it is not an evaluation of unknown-boundary recovery.

\subsection{Shared neural material with spatial descriptors}
The pretrained energy decoder is kept fixed. We infer a small spatial
material descriptor from the calibration prefix. For material coordinate
$X\in[0,1]^3$, the scale field is
\begin{equation}
 s_b(X)=s_0+\mathbf{1}_{b=2}\Delta s+
 c\tanh\!\left(8(X_1-r)\right),
\end{equation}
where $b$ indexes the deformable body. The same field parameterization is
available in every family; the heterogeneous-scene label does not switch
on a privileged material model. A second body receives an additional scale
offset, while the spatial profile and compressibility parameter are shared.
The other two pretrained descriptor coordinates are fixed at their
values $(0,-1)$ in this experiment.

We extend the frozen decoder with an independent positive multiplier on
its volumetric contribution. Writing $W_\theta$ for the pretrained energy,
$J=\det F$, and $B_\theta(s)$ for its existing positive volumetric
coefficient, the extension is
\begin{equation}
 W_{\theta,\beta}(F,s)=W_\theta(F,s)
 +(e^\beta-1)B_\theta(s)(J-\log J-1).
\end{equation}
Thus $\beta=0$ recovers the previous model exactly. The descriptor parameters
are fitted by bounded derivative-free optimization of the observed
point-position loss. No network weights are updated on a held world.
Additional calibration stages, when selected in development, initialize
from the preceding sealed best fit to the identical observation prefix.
This changes the optimization budget without reading the future query.
The frozen schedule permits $64$ objective evaluations followed by up to
$128$ additional evaluations, initialized from the best calibration fit.
The analytic control receives the identical schedule. The best calibration
loss determines the parameters used for both queries; no query-specific
parameter selection is performed. Reaching an evaluation limit is reported
separately from optimizer convergence.

Predicted motion follows the coupled mechanics rollout with explicit
collision geometry and compliant fixtures. The current scene model includes
conservative penalty contact, separation, and translation of the rigid
sphere; it does not model general rotating rigid shapes, friction, or
continuous collision detection. These limits define the evaluated scene
capabilities rather than being hidden in an aggregate score.

\subsection{Adapted mechanics baselines}
We compare with PhysTwin~\cite{jiang2025phystwin},
Spring-Gaus~\cite{zhong2024springgaus}, and
UniPhy~\cite{mittal2025uniphy}, using pinned author implementations.
PhysTwin uses the author Warp spring-force kernel and native forecast
integration, with per-edge stiffness fitted to the common track prefix.
Spring-Gaus uses the author internal forces and integration, a scalar
stiffness initialization search, and subsequent per-edge fitting. Both use
known object masses and explicit common scene-boundary adapters. PhysTwin
uses fixed zero dashpot and drag damping in this adapter. These are
mechanics comparisons with disclosed interfaces, not reproductions of the
complete author image-reconstruction pipelines.
Both native spring adapters use $40$ Adam iterations with learning rate
$0.03$ and a four-hour fitting limit. Spring-Gaus first searches a scalar
stiffness initialization. Single-clamp development controls at $160$
iterations improved calibration but increased vector forecast error for
both adapters; these controls and their shape errors are retained alongside
the selected budgets. This observation does not establish that $40$
iterations is optimal for every scene. Actual successful iterations and
numerical failures are reported rather than assumed from requested budgets.

UniPhy uses the released pretrained constitutive networks and native MPM,
with a fitted material latent and explicit bowl/obstacle hooks. Its current
adapter supports the drop family only. Its native separating-velocity
projection differs from the conservative penalty contact of the reference
and the other adapters; its results therefore measure this adapted
system's capability, not a matched-contact-law material-inference error.
Any adjoint adaptation and its derivative checks are reported separately.
Unsupported configurations are retained as unsupported, rather than being
counted as failures or wins. An analytic Neo-Hookean control uses the same
spatial fitting parameterization but has the advantage of knowing the
reference constitutive family.
UniPhy uses eight fixed-seed latent-centroid initializations followed by
eight AdamW iterations at learning rate $3\times10^{-5}$. Its derivative
qualification remains coordinate-sensitive, which limits attribution of
any adapted-system failure to the original method.

\subsection{Physical and visual evaluation}
We report full-query point-position RMS and per-object rigid-aligned shape
RMS separately. Additional physical diagnostics include sampled contact
penetration and first-contact time, net nongravitational body impulse,
and signed tetrahedron-volume ratios on the same material grid for all
methods. The impulse is computed from saved velocity and known mass after
subtracting gravity; it does not resolve individual contact forces.
Sampled orientation and gap checks are diagnostics, not proofs of global
injectivity or continuous nonintersection.

For a controlled image comparison, every method and the reference use the
same frozen Gaussian appearance, the same motion-transport map, and fixed
withheld cameras. We report PSNR, SSIM, LPIPS, and silhouette IoU, averaging
noninitial query frames for the primary scores and retaining all-frame
values separately.
This separates motion-induced image error from independent appearance
optimization. The rendering is object-only and does not test environmental
occlusion or native end-to-end reconstruction quality. Image agreement and
physical agreement are evaluated separately; neither is inferred from the
other.

Numerical reference uncertainty is assessed through spatial refinement,
with dense sampled orientation checks. A completed reference trajectory is
not automatically treated as spatially converged. Fitting, rollout,
reference and infrastructure failures remain separately visible in the
complete planned census. We average the two queries within each world,
then compute paired method differences and bootstrap intervals by resampling
worlds, not frames. Error comparisons conditional on successful paired
worlds are accompanied by their completion and coverage counts.
References use Q2 meshes with $(4n,n,n)$ cells, $n=4$, Lobatto-4 quadrature,
and a $25$\,$\mu$s time step. Both queries of the preselected first world
in every family are repeated at $n=6$. These twelve comparisons quantify
resolution sensitivity on those worlds; they are not rigorous error bounds
or a population-wide convergence certificate. Our material model and the
analytic control use $n=1$ with $50$\,$\mu$s steps. Native spring adapters
use $25$\,$\mu$s steps and UniPhy uses $50$\,$\mu$s steps.
